\subsection{Normalization and column gauge}
Let $B_0,B_1,B_2,B_3$ be unitary bases of $\mathbb C^6$.
After common left multiplication by $B_0^*$, put
$H_a=\sqrt6\,B_0^*B_a$ for $a=1,2,3$.
The quartet equations are
\[
 |(H_a)_{ij}|^2=1,\qquad H_a^*H_a=6I,
 \qquad |(H_a^*H_b)_{ij}|^2=6\quad(a\ne b).
\]
Fix $H_1$. Right multiplication of $H_2,H_3$ by independent diagonal
unitaries preserves these equations. Choosing diagonal entries
$\overline{(H_a)_{0j}}$ fixes both first rows to one.
The acting phase group is $\mathbb T^{12}$, of real dimension twelve.
The gauge fixes twenty-four ambient real coordinates and also discharges
twelve already-imposed circle equations. Thus $144-24=120$ ambient real
coordinates remain. This is not a claim about solution-set dimension,
smoothness, or Jacobian rank. If the dephased $H_1$ varies, its $5\times5$
block contributes fifty additional ambient coordinates.
Concrete normalization and simultaneous monomial transport were replayed in
\texttt{MubQuartetTransportBridge.lean} and
\texttt{MubMonomialTransport.lean}. Matching an executable coordinate
contractor to this gauge remains a separate obligation.

\subsection{Balance equations are ideal consequences}
Write $A=H_2$, $B=H_3$, $W=A^*B$, and define
\[
 g_{A,j}=A_{0j}-1,\qquad
 R_{B,m}=\sum_kB_{mk}\overline{B_{0k}}-6\delta_{m0}.
\]
For arbitrary complex matrices, before imposing any equations,
\begin{align*}
 \sum_kW_{jk}-6
 &=\sum_m\overline{A_{mj}}R_{B,m}
   -\sum_{m,k}\overline{A_{mj}}B_{mk}\overline{g_{B,k}}
   +6\overline{g_{A,j}},\\
 \sum_jW_{jk}-6
 &=\sum_m\overline{R_{A,m}}B_{mk}
   -\sum_{m,j}g_{A,j}\overline{A_{mj}}B_{mk}+6g_{B,k}.
\end{align*}
The first identity follows by substituting
$\overline{B_{0k}}=1+\overline{g_{B,k}}$; the second is its conjugate
with $A,B$ interchanged. Their real and imaginary parts lie in the ideal
generated by the corresponding gauge and row-Gram residuals.
Square-matrix column unitarity implies row unitarity, so all twelve complex
balance equations vanish on the quartet variety. Adding these consequences
does not restrict that variety further, although it may strengthen an
interval contractor. No independence claim for the other equations follows.
The two identities and ideal-membership corollaries were replayed over an
arbitrary commutative star ring in \texttt{MubGaugeSyzygy.lean}.

\subsection{Exact Tao evaluation and exclusion boundaries}
Let $\omega=-1/2+i\sqrt3/2$ and $S_{rj}=\omega^{E_{rj}}$, where
\[
 E=\begin{pmatrix}
 0&0&0&0&0&0\\0&0&1&1&2&2\\0&1&0&2&2&1\\
 0&1&2&0&1&2\\0&2&2&1&0&1\\0&2&1&2&1&0
 \end{pmatrix}.
\]
For every sign vector $\gamma\in\{-1,1\}^6$ with three positive entries,
\[
 \left|\sum_{j=0}^{5}\prod_{r=0}^{5}S_{rj}^{\gamma_r}\right|^2=9.
\]
All twenty masks were checked exactly; \texttt{MubTaoLaurent.lean}
connects the integer computation to complex multiplication and Laurent
powers. This evaluation does not establish a vanishing condition on
arbitrary MUB triples or independently certify Tao nonextension.

Jaming et al., arXiv:0902.0882v2, exclude a quartet containing the identity
and a member of their continuous Fourier family. Brierley and Weigert,
arXiv:0901.4051v2, Section 4.1, report the Tao common-vector obstruction.
Applying either result requires exact representative matching and the
appropriate common-unitary/monomial transport. The inherited unspecified
Haagerup ceiling is not admitted without concrete entries and a source theorem.
Maxwell and Brierley, arXiv:1402.4070v2, treat the full three-parameter
Karlsson family and a Laurent identity; vanishing alone gives no nonextension.
The inherited five-family completeness premise is quarantined.
A finite-corner classification concerns individual Hadamards and still
requires exclusion for every retained reconstruction over its continuous
seed domain. It does not by itself bound the number of mutually unbiased bases.

\paragraph{Claim scope.}
These results support concrete gauge transport, exact ideal consequences,
and the displayed Tao evaluation. They do not support $N(6)=3$, global
nonextension, a five-family cover, or a 120-dimensional solution manifold.
The narrowest claim is the stated matrix identities and exact evaluation.
The next decisive test is a quantified exclusion on the complete reconstruction
domain, or an exact quartet witness. This section is one agent's consolidation;
the five staged merge assignments are not represented as executed reviews.
